\section{Jump Game II}
\label{sec:jump-game-2}

\problemstatement{We are given an array $\textit{nums}$ of length $n$, where
$\textit{nums}[i]$ is the maximum number of steps we may jump forward from
index $i$. It is guaranteed that index $n-1$ is reachable from index $0$.
Starting at index $0$, we must find the \textbf{minimum number of jumps}
required to reach index $n-1$.}

\begin{patternbox}
This is the natural sequel to Section~\ref{sec:jump-game}: the keywords
change from \emph{``is it possible''} to \emph{``minimum number of
jumps,''} which upgrades the problem from a feasibility check to an
optimization. Whenever a reachability-horizon greedy problem is upgraded to
ask for the \emph{minimum number of steps/jumps/levels} to get somewhere,
that is a strong signal to think in terms of \emph{BFS levels} -- here,
each ``level'' is exactly one jump, and we want to know how many levels it
takes for the frontier of reachable positions to engulf the last index. We
will see that this BFS-by-levels idea can be implemented without an actual
queue, using two horizon variables instead.
\end{patternbox}

\subsection*{Understanding the problem carefully}

Think of the jumps as expanding in waves, exactly like breadth-first search
on a graph where node $i$ has an edge to every node in
$[i+1, i+\textit{nums}[i]]$. After $0$ jumps, we have only reached index $0$.
After $1$ jump, we have reached every index reachable directly from index
$0$ -- a contiguous range $[1, \textit{nums}[0]]$. After $2$ jumps, we have
reached every index reachable from \emph{any} index in that first range, and
so on. The number of jumps needed to first include index $n-1$ in the
reachable range is exactly the answer we want.

The key realization is that we do not need an actual BFS queue holding every
individual index. Instead, at each ``wave'' (each fixed number of jumps used
so far), we only need to know the boundary of the current range, and the
farthest index reachable by spending one more jump from anywhere inside the
current range.

\begin{ideabox}[Greedy Choice]
Maintain two horizons: $\textit{currentEnd}$, the farthest index reachable
using the jumps counted so far, and $\textit{farthest}$, the farthest index
reachable using \emph{one more} jump from anywhere within the current
range. Scan left to right; for every index $i$ within the current range,
update $\textit{farthest} \gets \max(\textit{farthest}, i +
\textit{nums}[i])$. The instant we finish processing the current range (when
$i$ reaches $\textit{currentEnd}$), we have exhausted everything reachable
with the current jump count, so we commit to spending one more jump:
increment the jump counter and set $\textit{currentEnd} \gets
\textit{farthest}$, beginning the next wave.
\end{ideabox}

\subsection*{Why this greedy expansion is correct}

This is exactly BFS, level by level, except that instead of storing every
node at the current level in a queue, we only ever need the single number
$\textit{farthest}$ -- the best possible reach from \emph{any} node at the
current level -- because, just as in Jump Game~I, every index between the
previous horizon and the new farthest reach is automatically included in
this level's frontier the moment any node within the previous range can
reach it. We never need to know \emph{which} specific index within the
current range produced the best $\textit{farthest}$ value, only the value
itself, since the next wave starts from the entire new range regardless of
which earlier index unlocked it. This is what allows an $O(n)$ single pass
to replace what would otherwise look like an explicit multi-source BFS.

\subsection*{Algorithm}

\begin{enumerate}
  \item Initialize $\textit{jumps} \gets 0$, $\textit{currentEnd} \gets 0$,
        $\textit{farthest} \gets 0$.
  \item For $i \gets 0$ to $n-2$ (we never need to jump \emph{from} the last
        index):
  \begin{enumerate}
    \item $\textit{farthest} \gets \max(\textit{farthest},\, i +
          \textit{nums}[i])$.
    \item If $i = \textit{currentEnd}$: this wave is exhausted, so spend one
          more jump. $\textit{jumps} \gets \textit{jumps} + 1$;
          $\textit{currentEnd} \gets \textit{farthest}$.
  \end{enumerate}
  \item Return \textit{jumps}.
\end{enumerate}

\begin{algorithm}[H]
\caption{Jump Game II (Minimum Jumps)}
\begin{algorithmic}[1]
\Procedure{MinJumps}{\textit{nums}}
  \State $n \gets \text{length}(\textit{nums})$
  \State $\textit{jumps} \gets 0$, $\textit{currentEnd} \gets 0$, $\textit{farthest} \gets 0$
  \For{$i \gets 0$ to $n-2$}
    \State $\textit{farthest} \gets \max(\textit{farthest}, i + \textit{nums}[i])$
    \If{$i = \textit{currentEnd}$}
      \State $\textit{jumps} \gets \textit{jumps} + 1$
      \State $\textit{currentEnd} \gets \textit{farthest}$
    \EndIf
  \EndFor
  \State \Return \textit{jumps}
\EndProcedure
\end{algorithmic}
\end{algorithm}

\subsection*{Dry run}

\dryrun{Take $\textit{nums} = [2, 3, 1, 1, 4]$ ($n=5$, indices $0..4$).}

Start with $\textit{jumps}=0$, $\textit{currentEnd}=0$,
$\textit{farthest}=0$. We loop $i$ from $0$ to $3$ (i.e.\ $n-2=3$).

\begin{itemize}
  \item $i=0$: $\textit{farthest} = \max(0, 0+2) = 2$. Since
        $i = 0 = \textit{currentEnd}$, spend a jump: $\textit{jumps}=1$,
        $\textit{currentEnd} = 2$.
  \item $i=1$: $\textit{farthest} = \max(2, 1+3) = 4$. Since
        $i = 1 \ne \textit{currentEnd}=2$, no new jump yet.
  \item $i=2$: $\textit{farthest} = \max(4, 2+1) = 4$. Since
        $i = 2 = \textit{currentEnd}$, spend a jump: $\textit{jumps}=2$,
        $\textit{currentEnd} = 4$.
  \item $i=3$: $\textit{farthest} = \max(4, 3+1) = 4$. Since
        $i = 3 \ne \textit{currentEnd}=4$, no new jump.
\end{itemize}

The loop ends (we stop before $i=4$, the last index). The answer is
$\textit{jumps} = 2$: jump from index $0$ to index $1$ (using $\textit{nums}[0]=2$
to reach as far as index $1$, the start of the best next wave), then jump
from index $1$ directly to index $4$ using $\textit{nums}[1]=3$.

\subsection*{C++ implementation}

\begin{lstlisting}
#include <bits/stdc++.h>
using namespace std;

int minJumps(vector<int>& nums) {
    int n = nums.size();
    int jumps = 0, currentEnd = 0, farthest = 0;

    for (int i = 0; i < n - 1; i++) {
        farthest = max(farthest, i + nums[i]);
        if (i == currentEnd) {
            jumps++;
            currentEnd = farthest;
        }
    }
    return jumps;
}

int main() {
    vector<int> nums = {2, 3, 1, 1, 4};
    cout << "Minimum jumps: " << minJumps(nums) << endl;
    return 0;
}
\end{lstlisting}

\begin{complexitybox}
\textbf{Time Complexity.} The single loop visits each of the first $n-1$
indices exactly once, doing constant work per index, giving
\[
  O(n).
\]
\textbf{Space Complexity.} We maintain only three scalar variables
($\textit{jumps}$, $\textit{currentEnd}$, $\textit{farthest}$), so the space
complexity is $O(1)$.
\end{complexitybox}

\begin{takeawaybox}
A feasibility-horizon greedy (Section~\ref{sec:jump-game}) upgrades cleanly
into a minimum-steps greedy by reframing it as implicit, level-by-level BFS:
track the boundary of the current wave (\textit{currentEnd}) and the best
reach achievable from anywhere inside that wave (\textit{farthest}); each
time the scan catches up to the current boundary, a new wave -- one more
jump -- begins. This pattern of ``BFS without an explicit queue, using two
horizon pointers'' generalizes to any problem phrased as ``minimum number of
steps to extend reach to a target.''
\end{takeawaybox}
